%% HMRC Paper - LaTeX Version
%% 
%% HOW TO EDIT:
%% 1. Go to https://www.overleaf.com (free account)
%% 2. Click "New Project" → "Upload Project"
%% 3. Upload this .tex file
%% 4. Edit directly in browser!
%%
%% Or use any LaTeX editor: TeXShop, TeXstudio, VS Code + LaTeX Workshop

\documentclass[sigconf, nonacm]{acmart}
%% Use [sigconf] for SIGIR/CIKM style
%% Change to \documentclass{article} for simple format

\usepackage{booktabs}
\usepackage{graphicx}
\usepackage{amsmath}
\usepackage{algorithm}
\usepackage{algorithmic}
\usepackage{xcolor}

%% Remove ACM copyright for draft
\setcopyright{none}
\settopmatter{printacmref=false}
\renewcommand\footnotetextcopyrightpermission[1]{}
\pagestyle{plain}

%% ============================================================================
%% TITLE AND AUTHORS - EDIT THIS SECTION
%% ============================================================================

\title{HMRC: Heterogeneous Multi-Representative Centroids for Improved Vector Database Routing}

\author{Your Name}
\affiliation{
  \institution{Your Institution}
  \city{Your City}
  \country{Your Country}
}
\email{your.email@example.com}

%% Add more authors like this:
% \author{Co-Author Name}
% \affiliation{
%   \institution{Their Institution}
% }
% \email{coauthor@example.com}

%% ============================================================================
%% ABSTRACT
%% ============================================================================

\begin{abstract}
Modern vector databases use segment-based architectures where vectors are partitioned into segments, each represented by a centroid for query routing. We identify a critical failure mode: \textbf{heterogeneous segments}---containing vectors from multiple semantic topics---are poorly represented by their mean centroid, causing routing failures. Through experiments on MS MARCO (100K passages), we show that routing recall drops from 100\% on coherent segments to 64.4\% on realistic heterogeneous segments (mixing 3 topics). We propose \textbf{HMRC (Heterogeneous Multi-Representative Centroids)}, which uses k-means clustering to generate multiple representative centroids per segment. HMRC improves routing recall by \textbf{26\%} (64.4\% $\rightarrow$ 81.2\%) on semi-structured segments with only 3$\times$ centroid storage overhead, while maintaining 100\% recall on coherent segments.
\end{abstract}

%% ============================================================================
%% KEYWORDS
%% ============================================================================

\keywords{vector databases, approximate nearest neighbor search, routing, embeddings, RAG}

\begin{document}
\maketitle

%% ============================================================================
%% 1. INTRODUCTION
%% ============================================================================

\section{Introduction}

\subsection{Background}

Vector databases are foundational to modern retrieval-augmented generation (RAG) systems. To scale to billions of vectors, these systems employ \textbf{segment-based architectures} (also called IVF or partition-based indexing):

\begin{enumerate}
    \item Vectors are partitioned into segments (clusters)
    \item Each segment is represented by a \textbf{centroid} (typically the mean)
    \item At query time, the query is compared to centroids to \textbf{route} to the most promising segments
    \item Only vectors in selected segments are searched exhaustively
\end{enumerate}

This routing step is critical: if the correct segment is not selected, the target vector cannot be found regardless of downstream search quality.

\subsection{The Heterogeneous Segment Problem}

We identify that segment heterogeneity causes \textbf{centroid degradation}:

\begin{definition}[Segment Heterogeneity]
A segment $S$ is heterogeneous if it contains vectors from multiple distinct semantic regions, i.e., $S = \bigcup_{i=1}^{m} S_i$ where $m > 1$ and $\text{sim}(S_i, S_j) < \tau$ for $i \neq j$.
\end{definition}

\textbf{Problem:} For a heterogeneous segment, the mean centroid $\bar{\mathbf{c}} = \frac{1}{|S|}\sum_{\mathbf{v} \in S} \mathbf{v}$ may lie in a region of embedding space that represents \textbf{none} of the constituent topics.

\subsection{When Do Heterogeneous Segments Occur?}

Heterogeneous segments are common in production systems:

\begin{itemize}
    \item \textbf{Time-based batching:} Documents ingested in the same time window share a segment regardless of topic
    \item \textbf{Source-based partitioning:} Documents from the same source/tenant grouped together
    \item \textbf{Incremental updates:} New documents appended to existing segments
    \item \textbf{Imperfect clustering:} Initial clustering may not perfectly separate topics
\end{itemize}

\subsection{Contributions}

Our contributions are:
\begin{enumerate}
    \item \textbf{Problem Identification:} We formally characterize and empirically validate the heterogeneous segment problem on real data (Section~\ref{sec:problem})
    \item \textbf{Solution:} We propose HMRC, using multiple representative centroids per segment (Section~\ref{sec:method})
    \item \textbf{Evaluation:} We demonstrate 26\% routing improvement on realistic heterogeneous segments (Section~\ref{sec:experiments})
\end{enumerate}

%% ============================================================================
%% 2. RELATED WORK
%% ============================================================================

\section{Related Work}

\subsection{Inverted File Indexing (IVF)}

FAISS~\cite{johnson2019billion} and similar systems use IVF for scalable search. Vectors are clustered, and queries are routed to the nearest cluster(s). Systems like Milvus~\cite{wang2021milvus} and Pinecone use similar architectures.

\subsection{Multi-Probe Search}

To compensate for routing errors, systems search multiple clusters (nprobe parameter). However, this has \textbf{linear cost increase}: $k\times$ nprobe = $k\times$ search cost. Multi-probe does not fix the underlying centroid representation problem.

\subsection{SOAR}

Google's ScaNN~\cite{guo2020accelerating} uses SOAR (Spilling with Orthogonally Augmented Residuals) to assign vectors to multiple clusters. Key difference: SOAR operates at \textbf{vector level} (2-3$\times$ vector storage), while HMRC operates at \textbf{segment level} ($k\times$ centroid storage, negligible).

\subsection{HNSW}

Graph-based methods like HNSW~\cite{malkov2018efficient} avoid centroids entirely but have higher memory overhead and are not segment-aware for distributed systems.

\textbf{Gap:} No prior work addresses \textbf{segment-level centroid representation} for heterogeneous segments.

%% ============================================================================
%% 3. PROBLEM ANALYSIS
%% ============================================================================

\section{Problem Analysis}
\label{sec:problem}

\subsection{Centroid Degradation}

Consider a segment $S = S_1 \cup S_2$ containing vectors from two topics with means $\boldsymbol{\mu}_1$ and $\boldsymbol{\mu}_2$. The segment mean is:

\begin{equation}
\bar{\mathbf{c}} = \frac{|S_1|}{|S|}\boldsymbol{\mu}_1 + \frac{|S_2|}{|S|}\boldsymbol{\mu}_2
\end{equation}

For a query $\mathbf{q}$ near topic 1 ($\boldsymbol{\mu}_1$):

\begin{equation}
d(\mathbf{q}, \bar{\mathbf{c}}) > d(\mathbf{q}, \boldsymbol{\mu}_1)
\end{equation}

The mean centroid is \textbf{biased away} from both topics, potentially causing routing failures.

\subsection{Empirical Validation}

We validate this problem on MS MARCO passages (100K vectors, 384 dimensions). Table~\ref{tab:problem} shows routing recall under different segmentation strategies.

\begin{table}[h]
\caption{Routing Recall@10 degrades with segment heterogeneity. Coherent segments are created via k-means clustering. Semi-$k$ segments mix $k$ random topics per segment.}
\label{tab:problem}
\begin{tabular}{lcc}
\toprule
\textbf{Segmentation} & \textbf{Recall@10} & \textbf{Drop from Coherent} \\
\midrule
Coherent & 100.0\% & --- \\
Semi-3 (3 topics) & 64.4\% & -35.6\% \\
Semi-5 (5 topics) & 49.1\% & -50.9\% \\
Random & 14.6\% & -85.4\% \\
\bottomrule
\end{tabular}
\end{table}

\textbf{Finding:} Routing recall drops by 35.6\% on realistic heterogeneous segments (Semi-3), validating the severity of the problem.

%% ============================================================================
%% 4. METHOD
%% ============================================================================

\section{Method: HMRC}
\label{sec:method}

\subsection{Problem Formulation}

Let $D = \{\mathbf{v}_1, ..., \mathbf{v}_n\}$ be a corpus of $n$ vectors partitioned into $m$ segments $\{S_1, ..., S_m\}$.

\textbf{Routing Task:} Given query $\mathbf{q}$, select top-$k$ segments most likely to contain relevant vectors.

\textbf{Standard Approach:} Represent each segment by mean centroid:
\begin{equation}
\mathbf{c}_j = \frac{1}{|S_j|} \sum_{\mathbf{v} \in S_j} \mathbf{v}
\end{equation}

Route to segments with minimum distance:
\begin{equation}
\text{Route}(\mathbf{q}, k) = \underset{j \in [m]}{\text{argmin-}k} \; d(\mathbf{q}, \mathbf{c}_j)
\end{equation}

\subsection{HMRC: Multiple Representatives}

\textbf{Key Idea:} Represent each segment $S_j$ with \textbf{multiple} centroids $\{\mathbf{c}_j^{(1)}, ..., \mathbf{c}_j^{(r)}\}$ that capture its internal structure.

\begin{algorithm}[h]
\caption{HMRC Index Construction}
\label{alg:build}
\begin{algorithmic}[1]
\REQUIRE Segments $\{S_1, ..., S_m\}$, representatives $r$
\ENSURE Representative set $R$, segment mapping $M$
\FOR{$j = 1$ to $m$}
    \IF{$|S_j| < 2r$}
        \STATE $R_j \leftarrow \{\text{mean}(S_j)\}$ \COMMENT{Fallback}
    \ELSE
        \STATE $R_j \leftarrow \text{KMeans}(S_j, k=r).\text{centroids}$
    \ENDIF
    \STATE $R \leftarrow R \cup R_j$
    \FOR{each $\mathbf{c} \in R_j$}
        \STATE $M[\mathbf{c}] \leftarrow j$ \COMMENT{Map to segment}
    \ENDFOR
\ENDFOR
\RETURN $R$, $M$
\end{algorithmic}
\end{algorithm}

\begin{algorithm}[h]
\caption{HMRC Query Routing}
\label{alg:route}
\begin{algorithmic}[1]
\REQUIRE Query $\mathbf{q}$, representatives $R$, mapping $M$, $k$
\ENSURE Top-$k$ segment IDs
\STATE $\text{nearest} \leftarrow \text{argmin}_{\mathbf{c} \in R} d(\mathbf{q}, \mathbf{c})$ \COMMENT{Top $2rk$}
\STATE $\text{segments} \leftarrow []$
\FOR{each $\mathbf{c} \in \text{nearest}$}
    \STATE $s \leftarrow M[\mathbf{c}]$
    \IF{$s \notin \text{segments}$}
        \STATE $\text{segments}.\text{append}(s)$
    \ENDIF
    \IF{$|\text{segments}| \geq k$}
        \STATE \textbf{break}
    \ENDIF
\ENDFOR
\RETURN $\text{segments}[:k]$
\end{algorithmic}
\end{algorithm}

\subsection{Complexity Analysis}

\textbf{Storage Overhead:}
\begin{itemize}
    \item Baseline: $m$ centroids $\times$ $d$ dimensions = $O(md)$
    \item HMRC: $rm$ centroids $\times$ $d$ dimensions = $O(rmd)$
    \item \textbf{Overhead factor:} $r$ (typically 3-5$\times$)
\end{itemize}

For a system with 10K segments and 384-dim vectors:
\begin{itemize}
    \item Baseline: $10K \times 384 \times 4$ bytes = \textbf{15 MB}
    \item HMRC-3: $30K \times 384 \times 4$ bytes = \textbf{45 MB}
\end{itemize}

This is \textbf{negligible} compared to vector storage (GBs to TBs).

%% ============================================================================
%% 5. EXPERIMENTS
%% ============================================================================

\section{Experiments}
\label{sec:experiments}

\subsection{Experimental Setup}

\textbf{Dataset:} MS MARCO passages, 100,000 passages embedded with all-MiniLM-L6-v2 (384 dimensions).

\textbf{Segmentation:} 
\begin{itemize}
    \item \textbf{Coherent:} K-means clustering (200 segments)
    \item \textbf{Semi-$k$:} Each segment mixes $k$ random topic clusters
    \item \textbf{Random:} Uniform random assignment
\end{itemize}

\textbf{Queries:} 2,000 random passages used as queries; ground truth is the segment containing each query vector.

\textbf{Metric:} Routing Recall@$k$ --- fraction of queries where the ground-truth segment is in the top-$k$ routed segments:
\begin{equation}
\text{Recall@}k = \frac{1}{|Q|} \sum_{\mathbf{q} \in Q} \mathbb{1}[s^*(\mathbf{q}) \in \text{Route}(\mathbf{q}, k)]
\end{equation}

\subsection{Main Results}

Table~\ref{tab:main} shows HMRC performance on Semi-3 segments (realistic heterogeneity).

\begin{table}[h]
\caption{HMRC improves routing recall on semi-structured segments. Results on Semi-3 segmentation (3 topics per segment).}
\label{tab:main}
\begin{tabular}{lccc}
\toprule
\textbf{Method} & \textbf{Recall@10} & \textbf{Improvement} & \textbf{Overhead} \\
\midrule
Baseline & 64.4\% & --- & 1.0$\times$ \\
HMRC-2 & 72.1\% & +12\% & 2.0$\times$ \\
\textbf{HMRC-3} & \textbf{81.2\%} & \textbf{+26\%} & \textbf{3.0$\times$} \\
HMRC-5 & 80.7\% & +25\% & 5.0$\times$ \\
HMRC-7 & 80.8\% & +26\% & 7.0$\times$ \\
\bottomrule
\end{tabular}
\end{table}

\textbf{Finding:} HMRC-3 achieves optimal performance with 26\% improvement at only 3$\times$ centroid storage overhead.

\subsection{No Harm on Coherent Segments}

Table~\ref{tab:noharm} shows HMRC maintains performance on coherent segments.

\begin{table}[h]
\caption{HMRC does not harm coherent segment performance.}
\label{tab:noharm}
\begin{tabular}{lccc}
\toprule
\textbf{Method} & \textbf{Coherent} & \textbf{Semi-3} & \textbf{Random} \\
\midrule
Baseline & 100.0\% & 64.4\% & 14.6\% \\
HMRC-3 & 100.0\% & 81.2\% & 16.2\% \\
HMRC-5 & 100.0\% & 80.7\% & 22.1\% \\
\bottomrule
\end{tabular}
\end{table}

\textbf{Finding:} HMRC maintains 100\% recall on coherent segments while providing substantial gains on heterogeneous segments.

\subsection{Ablation: Number of Representatives}

Figure~\ref{fig:ablation} shows the effect of varying the number of representatives $r$.

%% UNCOMMENT WHEN YOU HAVE FIGURES:
% \begin{figure}[h]
%     \centering
%     \includegraphics[width=0.8\columnwidth]{fig4_representatives_ablation.png}
%     \caption{Recall vs. number of representatives. Diminishing returns after $r=3$.}
%     \label{fig:ablation}
% \end{figure}

\begin{table}[h]
\caption{Ablation on number of representatives (Semi-3 segments).}
\label{tab:ablation}
\begin{tabular}{ccc}
\toprule
\textbf{Representatives ($r$)} & \textbf{Recall@10} & \textbf{Overhead} \\
\midrule
1 (Baseline) & 64.4\% & 1$\times$ \\
2 & 72.1\% & 2$\times$ \\
3 & 81.2\% & 3$\times$ \\
4 & 81.0\% & 4$\times$ \\
5 & 80.7\% & 5$\times$ \\
7 & 80.8\% & 7$\times$ \\
10 & 81.5\% & 10$\times$ \\
\bottomrule
\end{tabular}
\end{table}

\textbf{Finding:} Performance plateaus at $r=3$, making it the optimal choice for Semi-3 segments.

\subsection{Ablation: Heterogeneity Level}

Table~\ref{tab:topics} shows HMRC benefit increases with segment heterogeneity.

\begin{table}[h]
\caption{HMRC benefit increases with heterogeneity level.}
\label{tab:topics}
\begin{tabular}{lccc}
\toprule
\textbf{Topics/Segment} & \textbf{Baseline} & \textbf{HMRC-3} & \textbf{Relative Gain} \\
\midrule
1 (Coherent) & 100.0\% & 100.0\% & 0\% \\
2 & 73.2\% & 80.1\% & +9\% \\
3 & 64.4\% & 81.2\% & +26\% \\
5 & 49.1\% & 63.0\% & +28\% \\
7 & 40.2\% & 49.6\% & +23\% \\
10 & 30.5\% & 41.0\% & +34\% \\
\bottomrule
\end{tabular}
\end{table}

\textbf{Finding:} HMRC provides consistent improvements across all heterogeneity levels.

\subsection{Ablation: Segment Size}

Table~\ref{tab:sizes} shows HMRC works across different segment sizes.

\begin{table}[h]
\caption{HMRC improvement across segment sizes (Semi-3 segments).}
\label{tab:sizes}
\begin{tabular}{lccc}
\toprule
\textbf{Segment Size} & \textbf{Baseline} & \textbf{HMRC-3} & \textbf{Improvement} \\
\midrule
100 & 36.0\% & 72.0\% & +100\% \\
250 & 50.0\% & 77.0\% & +54\% \\
500 & 64.4\% & 81.2\% & +26\% \\
1000 & 69.0\% & 77.0\% & +12\% \\
2000 & 72.0\% & 75.0\% & +4\% \\
\bottomrule
\end{tabular}
\end{table}

\textbf{Finding:} HMRC provides larger improvements for smaller segments, which are more heterogeneous relative to their size.

%% ============================================================================
%% 6. DISCUSSION
%% ============================================================================

\section{Discussion}

\subsection{When Does HMRC Help?}

HMRC provides significant improvements when:
\begin{enumerate}
    \item Segments contain \textbf{discoverable multimodal structure}
    \item Internal topics are separable by k-means
    \item Heterogeneity level roughly matches representative count
\end{enumerate}

HMRC provides minimal improvement on \textbf{uniformly random} segments (no structure to discover) and \textbf{perfectly coherent} segments (mean is already optimal).

\subsection{Practical Recommendations}

Based on our experiments:
\begin{itemize}
    \item \textbf{Default:} $r=3$ representatives (good for most workloads)
    \item \textbf{High heterogeneity:} $r=5$ for highly mixed segments
    \item \textbf{Memory constrained:} $r=2$ still provides $\sim$12\% improvement
\end{itemize}

\subsection{Integration}

HMRC is a \textbf{drop-in replacement} for existing IVF systems:
\begin{itemize}
    \item Same query interface
    \item Same segment structure
    \item Only centroid computation changes
\end{itemize}

%% ============================================================================
%% 7. CONCLUSION
%% ============================================================================

\section{Conclusion}

We identified the \textbf{heterogeneous segment problem} in vector database routing and proposed \textbf{HMRC} as a simple, effective solution. Key results:

\begin{itemize}
    \item \textbf{Problem:} 35.6\% routing recall drop on realistic heterogeneous segments
    \item \textbf{Solution:} HMRC with k-means clustering per segment
    \item \textbf{Improvement:} +26\% recall with 3$\times$ centroid storage overhead
    \item \textbf{No harm:} Maintains 100\% recall on coherent segments
\end{itemize}

HMRC addresses a fundamental gap in vector database routing without requiring changes to downstream search or application logic.

%% ============================================================================
%% REFERENCES
%% ============================================================================

\bibliographystyle{ACM-Reference-Format}

\begin{thebibliography}{9}

\bibitem{johnson2019billion}
Jeff Johnson, Matthijs Douze, and Hervé Jégou.
\newblock Billion-scale similarity search with GPUs.
\newblock {\em IEEE Transactions on Big Data}, 7(3):535--547, 2019.

\bibitem{guo2020accelerating}
Ruiqi Guo, Philip Sun, Erik Lindgren, Quan Geng, David Simcha, Felix Chern, and Sanjiv Kumar.
\newblock Accelerating large-scale inference with anisotropic vector quantization.
\newblock In {\em International Conference on Machine Learning}, pages 3887--3896, 2020.

\bibitem{malkov2018efficient}
Yu A. Malkov and Dmitry A. Yashunin.
\newblock Efficient and robust approximate nearest neighbor search using hierarchical navigable small world graphs.
\newblock {\em IEEE Transactions on Pattern Analysis and Machine Intelligence}, 42(4):824--836, 2018.

\bibitem{wang2021milvus}
Jianguo Wang, Xiaomeng Yi, Rentong Guo, Hai Jin, Peng Xu, Shengjun Li, Xiangyu Wang, Xiangzhou Guo, Chengming Li, Xiaohai Xu, et al.
\newblock Milvus: A purpose-built vector data management system.
\newblock In {\em Proceedings of the 2021 International Conference on Management of Data}, pages 2614--2627, 2021.

\bibitem{lewis2020retrieval}
Patrick Lewis, Ethan Perez, Aleksandra Piktus, Fabio Petroni, Vladimir Karpukhin, Naman Goyal, Heinrich Küttler, Mike Lewis, Wen-tau Yih, Tim Rocktäschel, et al.
\newblock Retrieval-augmented generation for knowledge-intensive NLP tasks.
\newblock {\em Advances in Neural Information Processing Systems}, 33:9459--9474, 2020.

\end{thebibliography}

\end{document}
